% =============================================================================
%  II. BACKGROUND: LIMITATIONS OF THE CONVENTIONAL APPROACH
% =============================================================================
\begin{frame}{II. Background: Limitations of the Conventional Approach}
\vspace*{0pt}%
{\centering One label $\rightarrow$ one split $\rightarrow$ AUROC
$\rightarrow$ model comparison\par}

\vspace{2pt}
\begin{enumerate}\itemsep1pt\small
  \item \textbf{Label definition:} applying Sepsis-3 differently produced
        867--2,178 cases in the same database~\cite{cohen2024subtle}.
  \item \textbf{Data split:} sepsis prediction showed the largest performance
        loss under temporal shift among the clinical tasks
        evaluated~\cite{guo2022temporal}.
  \item \textbf{Metric choice:} high AUROC can coexist with utility below the
        no-alert strategy~\cite{wang2025srpm}.
  \item \textbf{Clinical-action dependence:} labels may capture recognition and
        treatment practice, not physiological deterioration
        alone~\cite{kamran2024evaluation}.
\end{enumerate}

\vspace{2pt}
\begin{block}{This thesis}
  \small Evaluates model class, label definition, data split, prediction
  timing and evaluation metric as separate \textbf{sources of variation}, with
  \textbf{bootstrap confidence intervals} so that their magnitudes can be
  compared.
\end{block}
\end{frame}
